% SPDX-License-Identifier: Apache-2.0
\section{A Lattice of Nodes}
\label{sec:x-mesh}

\code{Mesh<W, H, N, ..>} from the network on chip~\cite{noc} is the
largest array of children in the tree: \code{N} nodes, \code{W} wide
and \code{H} high, in one field of \code{Units} (issue 635). Every node
is one type, because a node's column and row are two inputs rather
than two type parameters, and the mesh holds them at constants with
\code{tie}. The index may appear in the constant, so node \code{i} is
given \code{tie(U::<XB>::from(i \% W))} and
\code{tie(U::<YB>::from(i / W))}, and the lowering gives each node
wires of its own driven by those values.

The links are one set of channels per direction of travel and per
virtual channel, made with \code{chans}. Node \code{i} sends up on
channel \code{i} of the upward set, and the node above reads that
channel on its south input. A link at an edge wraps to the other end
of its row or column, so that every channel has one driver and one
reader. Routing is dimension order, so a packet for a place on the
lattice never takes one of those links.

The bench makes a mesh three wide and two high. Every node sends a
request to each of the other five, and a response to the node at the
opposite corner. The run asserts that each packet leaves by the exit
of the node it was addressed to, and that all thirty requests and six
responses arrive. The netlist, six node instances and the channels
between them, is checked against the run under nvc and Verilator.

\begin{figure}[htbp]
\centering
\input{mesh_timing.tex}
\caption{Six nodes: the requests offered at each exit, and where they
leave.}
\label{fig:x-mesh}
\end{figure}

\lstinputlisting[caption={\code{ex\_mesh.rs}. Run with \code{bazel run //lib/examples:ex\_mesh}.},label={lst:x-mesh}]{lib/examples/ex_mesh.rs}

\lstinputlisting[style=ascii,caption={What \code{ex\_mesh} printed when the build ran it: each packet as it left, then the instances.},label={lst:x-mesh-out}]{out_mesh.txt}
